\documentclass[12pt,a4paper]{report}

% -------------------------------------------------
% PACKAGES
% -------------------------------------------------
\usepackage[utf8]{inputenc}
\usepackage[T1]{fontenc}
\usepackage{lmodern}
\usepackage[english]{babel}
\usepackage{amsmath, amssymb}
\usepackage{graphicx}
\usepackage{float}
\usepackage{caption}
\captionsetup{font=small, labelfont=bf}
\captionsetup[figure]{position=bottom}
\captionsetup[table]{position=top}
\usepackage{subcaption}
\usepackage{geometry}
\usepackage{booktabs}
\usepackage{algorithm}
\usepackage{algpseudocode}
\usepackage{setspace}
\usepackage{titlesec}
\usepackage{enumitem}
\usepackage{tocloft}
\usepackage{hyperref}
\usepackage[nameinlink, noabbrev]{cleveref}

% -------------------------------------------------
% PAGE FORMAT
% -------------------------------------------------
\geometry{top=2.5cm, bottom=2.5cm, left=2.5cm, right=2.5cm}
\setstretch{1.2}
\setlength{\parskip}{6pt}
\setlength{\parindent}{0pt}
\raggedbottom
\setlength{\emergencystretch}{3em}
\sloppy
\setlist[itemize]{noitemsep, topsep=3pt}
\setlist[enumerate]{noitemsep, topsep=3pt}
\graphicspath{{figures/}}

\hypersetup{
  colorlinks=true,
  linkcolor=black,
  urlcolor=cyan,
  citecolor=blue
}

% Algorithms
\makeatletter
\setlength{\ALG@thistlm}{\leftmargin}
\makeatother

% -------------------------------------------------
% TITLE FORMAT
% -------------------------------------------------
\titleformat{\chapter}[hang]
  {\normalfont\huge\bfseries}
  {\thechapter.}{1em}{}
\titlespacing*{\chapter}{0pt}{-10pt}{10pt}

\titleformat{\section}
  {\normalfont\bfseries\large}{\thesection.}{0.5em}{}

\titleformat{\subsection}
  {\normalfont\bfseries\normalsize}{\thesubsection.}{0.5em}{}

% -------------------------------------------------
% TABLE OF CONTENTS
% -------------------------------------------------
\renewcommand{\cftbeforetoctitleskip}{-1em}
\renewcommand{\cftaftertoctitleskip}{1em}
\setlength{\cftchapindent}{0pt}
\setlength{\cftsecindent}{1.5em}
\setlength{\cftsubsecindent}{3em}
\setlength{\cftsubsubsecindent}{4.5em}
\renewcommand{\cftdotsep}{1}

\begin{document}

% -------------------------------------------------
% COVER PAGE
% -------------------------------------------------
\begin{titlepage}
    \centering
    \includegraphics[width=0.35\textwidth]{logo_ulpgc.png} \\[0.8cm]

    {\huge \textbf{Big Data}} \\[0.3cm]
    {\Large Degree in Data Science and Engineering} \\[0.2cm]
    {\large School of Computer Engineering} \\[0.2cm]
    {\large Academic year 2026--2027} \\[1cm]

    \line(1,0){400} \\[4mm]
    {\Huge \textbf{Individual Assignment 1}} \\[3mm]
    {\LARGE Basic Matrix Multiplication in Python, Java and C} \\[3mm]
    {\large \textit{How matrix size affects the runtime and memory use\\of a basic triple-loop matrix multiplication}} \\[4mm]
    \line(1,0){400} \\[1cm]

    {\large
    \begin{tabular}{l}
        \textbf{Name} \\
        \midrule
        Claudia Álvarez González \\
    \end{tabular}} \\[1cm]

    {\large \textbf{Repository:} \url{https://github.com/Claudia1771/bigdata-matrix-multiplication}} \\[2mm]
    {\normalsize (folder \texttt{assignment1})}

    \vfill

    {\large \today}
\end{titlepage}

% -------------------------------------------------
% INDEXES
% -------------------------------------------------
\pagenumbering{roman}
\tableofcontents
\newpage
\listoffigures
\newpage
\listoftables
\newpage
\listofalgorithms
\newpage
\pagenumbering{arabic}

%ABSTRACT%

% =================================================
\chapter{Introduction and Objective}
% =================================================

Dense matrix multiplication is a basic building block of scientific computing, computer graphics and machine
learning. For two $n \times n$ matrices the straightforward algorithm performs $n^3$ multiply-add operations, so its
cost grows cubically with the matrix side $n$. Asymptotic complexity, however, says nothing about the constant
factors that determine the actual running time on a given machine. These depend on how the language executes the code
(interpretation, just-in-time or ahead-of-time compilation), how the matrices are stored in memory and how the
algorithm interacts with the memory hierarchy~\cite{leiserson2020,drepper2007}.

\section{Research question}

\textit{How does the matrix size $n$ affect the execution time and the memory use of the same basic triple-loop
matrix multiplication implemented in Python, Java and C, under fixed and documented execution conditions?}

The factors changed are the matrix size $n$ (the side of an $n \times n$ square matrix) and the implementation
language. The outcomes measured are the execution time of one multiplication and the peak memory of the process.

\section{Hypotheses}

The following predictions were stated before running the experiment:
\begin{enumerate}[label=\textbf{H\arabic*.}]
  \item For each language the execution time grows approximately as $n^3$, so doubling $n$ multiplies the time by about 8.
  \item Python is one to two orders of magnitude slower than Java and C, because the interpreter executes every operation on boxed objects.
  \item Java and C are within a small factor of each other once the JIT compiler has optimised the loop.
  \item Memory grows as $n^2$, with a fixed runtime overhead that is largest for Java and a per-element overhead that is largest for Python.
\end{enumerate}

% =================================================
\chapter{Background}
% =================================================

\section{The basic algorithm}

Given $A$ and $B$, each element of $C = AB$ is $C_{ij} = \sum_{k=1}^{n} A_{ik} B_{kj}$~\cite{golub2013}.
\Cref{alg:triple} shows the classic $i$-$j$-$k$ triple loop used in this study. It performs exactly $n^3$
multiplications and $n^3$ additions.

\begin{algorithm}[H]
\caption{Basic $i$-$j$-$k$ matrix multiplication}
\label{alg:triple}
\begin{algorithmic}[1]
\Require $A, B \in \mathbb{R}^{n \times n}$
\Ensure $C = AB$
\State allocate $C \in \mathbb{R}^{n \times n}$
\For{$i \gets 1$ \textbf{to} $n$}
  \For{$j \gets 1$ \textbf{to} $n$}
    \State $s \gets 0$
    \For{$k \gets 1$ \textbf{to} $n$}
      \State $s \gets s + A_{ik} \cdot B_{kj}$
    \EndFor
    \State $C_{ij} \gets s$
  \EndFor
\EndFor
\State \Return $C$
\end{algorithmic}
\end{algorithm}

\section{Memory access pattern}

In the $i$-$j$-$k$ order the innermost loop reads one row of $A$ sequentially and one column of $B$ with a stride of
$n$ elements. With row-major storage, each access to $B$ touches a different cache line, so for large $n$ most
accesses to $B$ miss in the cache~\cite{drepper2007,hennessy2017,lam1991}. When the row length is a power of two, the
elements of a column map to a small number of cache sets and can evict each other (conflict
misses)~\cite{drepper2007,lam1991}. Optimised libraries avoid these effects with blocking and
vectorisation~\cite{goto2008}; they are outside the scope of this controlled comparison and will be studied in
Assignment~2.

\section{Execution models}

CPython interprets bytecode; every \texttt{float} is a heap object, and every arithmetic operation involves dynamic
type dispatch and reference counting. Java code is first interpreted and then compiled at run time by the HotSpot JIT
compiler, which is why warm-up runs are recommended before measuring~\cite{georges2007}. C is compiled ahead of time
to native code. Leiserson et al.~\cite{leiserson2020} used this same triple loop to show differences of several orders
of magnitude between Python and C.

\section{Floating-point arithmetic}

Floating-point addition is not associative, so a different summation order can change the last bits of a result. For
an inner product of length $n$, the error with respect to the exact value is bounded by about
$n u \sum_k |A_{ik} B_{kj}|$, where $u \approx 1.1 \times 10^{-16}$ is the unit roundoff of
\texttt{float64}~\cite{higham2002}. Since all three implementations perform the same operations in the same order,
and Java uses strict IEEE~754 semantics by default since version~17~\cite{jep306}, they are expected to produce
bit-identical results.

% =================================================
\chapter{Methodology}
% =================================================

\section{Implementations}

The three programs implement \Cref{alg:triple} with a scalar accumulator, use 64-bit floating-point numbers and do not
call any matrix library. They differ only in the unavoidable, idiomatic storage layout of each language
(\Cref{tab:impl}). All implementations are single-threaded. The C program creates no additional threads; the JVM
creates internal threads for JIT compilation and garbage collection, which in this study share the single core
assigned to the process.

\begin{table}[H]
\centering
\small
\caption{Implementations, storage layouts and fixed execution settings.}
\label{tab:impl}
\begin{tabular}{p{1.4cm}p{6.0cm}p{6.0cm}}
\toprule
\textbf{Language} & \textbf{Storage of each matrix} & \textbf{Execution settings (fixed for the whole study)} \\
\midrule
Python & List of $n$ lists of Python \texttt{float} objects (each element is a pointer to a 24-byte object) & CPython 3.13.0, default interpreter (experimental JIT not enabled) \\
\addlinespace
Java & \texttt{double[][]}: an array of $n$ independent row arrays of primitive doubles & JDK 21.0.5 HotSpot, default tiered JIT, serial collector, \texttt{-Xmx4g} \\
\addlinespace
C & One contiguous row-major \texttt{double} array of $n^2$ elements & GCC 16.2.0 (MinGW-w64), \texttt{-O2 -std=c11}, no \texttt{-march} flag \\
\bottomrule
\end{tabular}
\end{table}

\section{Inputs}

To guarantee that all languages multiply the same logical matrices, the inputs are generated once by
\texttt{generate\_inputs.py} and read by the three programs. For each size, $A$ and $B$ are drawn with NumPy's
\texttt{default\_rng} (PCG64) using seed $42 + n$, uniform values in $[0, 1)$, type \texttt{float64}, and stored as raw
little-endian row-major binary files. A manifest records the seed and the SHA-256 hash of every file.

The declared size sequence is $n \in \{64, 128, 256, 512, 1000, 1024, 2048\}$: powers of two, so that consecutive
sizes should differ by a factor of 8 in time, plus $n = 1000$ next to $1024$ to check whether the power-of-two size
has an effect of its own.

\section{Correctness checks}

Before benchmarking, \texttt{verify.py} runs every implementation on:
\begin{itemize}
  \item $1\times1$ and $2\times2$ matrices with known integer products (exact comparison);
  \item $AI = A$ and $A\,0 = 0$ for a random $8\times8$ matrix $A$ (exact comparison);
  \item random matrices with $n \in \{1, 2, 3, 17, 64, 100\}$, compared with \texttt{numpy.matmul} as the trusted
  reference using $\mathrm{rtol} = \mathrm{atol} = 10^{-12}$.
\end{itemize}
This tolerance is about 100 times larger than the error bound $n u$ for $n = 100$ and values in $[0,1)$ (no
cancellation is possible because all products are non-negative), but small enough to detect any indexing error. The
outputs of the three languages are also compared with each other bit by bit.

In addition, during the benchmark every multiplication returns a checksum: the sequential sum of all elements of $C$,
reported as the raw bits of the double. This guarantees that the result is consumed, so the compiler cannot remove the
computation, and allows checking that every repetition in every language produced the same result.

\section{Measurement protocol}

\begin{itemize}
  \item \textbf{Timer boundary (kernel measurement).} The timer covers only the multiplication function, including the
  allocation of the output matrix $C$ (an $O(n^2)$ cost, negligible compared with the $O(n^3)$ loop). Reading the
  input files, process start-up and the checksum are outside the timed region. Monotonic timers are used:
  \texttt{time.perf\_counter} (Python), \texttt{System.nanoTime} (Java) and \texttt{clock\_gettime(CLOCK\_MONOTONIC)} (C).
  \item \textbf{Processes, warm-up and repetitions.} Each (language, $n$) configuration runs in a fresh process that
  loads the inputs, executes the warm-up multiplications (not reported) and then the measured repetitions. Default:
  5 measured repetitions, with 3 warm-ups for Java (JIT compilation) and 1 for Python and C (caches and memory
  allocation). Costly configurations, defined as those whose predicted time (median of the previous size multiplied by
  $(n/n_{\text{prev}})^3$) exceeds 20~s, use 3 repetitions and 0 warm-ups (1 for Java), to keep the study feasible.
  For runs of minutes, the JIT compiles the loop during the first fraction of a second, so its effect is negligible.
  \item \textbf{CPU affinity and priority.} The processor is hybrid, and a probe that runs the same CPU-bound loop
  pinned to each logical processor (\texttt{core\_probe.py}, \Cref{sec:hybrid}) shows that the efficient cores are
  about 2--2.7 times slower than the performance cores. A preliminary run without affinity showed repetitions up to
  3.4 times slower than others of the same configuration. Therefore, every benchmark process is created already pinned
  to logical processor~12 (a performance core) with \texttt{HIGH} priority class. Because the JVM detects a single
  available processor, the garbage collector is fixed explicitly to the serial collector (\texttt{-XX:+UseSerialGC}),
  and JIT compilation shares the same core as the measured thread.
  \item \textbf{Budgets.} Maximum time of one multiplication: 300~s. Memory budget: 4~GiB per process. A configuration
  that exceeds a limit is stopped and recorded as a timeout or memory failure, and the larger sizes of that language
  are not attempted. No time is invented for failed configurations.
  \item \textbf{Summary statistics.} All raw measurements are kept in \texttt{raw\_measurements.csv}. Tables report the
  median, the first and third quartiles ($Q_1$, $Q_3$) and the interquartile range (IQR), computed with linear
  interpolation.
  \item \textbf{Memory metric.} Peak working set (peak resident memory) of the whole process, read with the Win32
  \texttt{GetProcessMemoryInfo} API after the process has finished. It includes the inputs, the result, temporary
  objects and the language runtime (interpreter or JVM), so it is not a measure of the matrix storage alone. To avoid
  measuring a launcher process, Java is started with the \texttt{java.exe} of the JDK. As a reference, the storage of
  three $n \times n$ \texttt{float64} matrices is $24n^2$ bytes.
\end{itemize}

\section{Environment}

\begin{table}[H]
\centering
\small
\caption{Hardware and software environment of the reported run.}
\label{tab:env}
\begin{tabular}{ll}
\toprule
\textbf{Item} & \textbf{Value} \\
\midrule
CPU & Intel Core Ultra 7 155H (hybrid: 6 P-cores with 2 threads each, 8 E-cores, 2 LP E-cores) \\
Cores / logical processors & 16 / 22 \\
Cache & 24 MB L3 \\
RAM & 16 GB \\
Operating system & Windows 11 (build 26200), laptop connected to the power supply \\
Python & CPython 3.13.0 (NumPy 2.2.6 only for inputs and reference) \\
Java & Oracle JDK 21.0.5 HotSpot, \texttt{-Xmx4g -XX:+UseSerialGC} \\
C compiler & GCC 16.2.0 (MinGW-w64 UCRT), \texttt{-O2 -std=c11} \\
CPU affinity & Logical processor 12 (performance core), \texttt{HIGH} priority class \\
\bottomrule
\end{tabular}
\end{table}

The whole configuration is stored in \texttt{config.json} and the environment captured at run time in
\texttt{environment.json}.

%RESULTS%

\end{document}
